% FIGURE 3 --- the capabilities figure. Unifying claim: ONE FROZEN R_theta
% supports multiple downstream design modes. Not a grab-bag of four results.
%
% ===================== DRAWING SPEC. READ BEFORE DRAWING. =====================
%
% LAYOUT: 2x2, escalating from standard to richer capability:
%   (a) standard editing -> (b) multi-objective control -> (c) state reuse after
%   a goal change -> (d) support-enforced admissibility.
%
% (a) SOURCE-CONDITIONED EDITING. A simple empirical curve.
%     x  candidate allowance k, over the MEASURED RANGE k = 1..12 ONLY
%     y  sources solved (of 128), or success %
%     cumulative solved: 31,40,44,48,51,56,61,63,65,67,68,70
%     Endpoint label 70/128 = 54.7% at k=12. (Confirmed against
%     docs/VALID128_DIVERSITY.json, n_sources_averaged = 70.)
%     Points connected by line segments.
%     Comparators as HORIZONTAL DASHED MARKERS at their reported success level,
%     not a duplicated table, and not a full comparator curve -- we do not have
%     their success as a function of k, so drawing one would be invented.
%     DO NOT draw beyond k=12. DO NOT put the pending 800-source extension in
%     the panel. DO NOT add the full-vs-size-fixed-support inset; it is
%     auxiliary and would clutter.
%
% (b) PREFERENCE-CONDITIONED MULTI-OBJECTIVE DESIGN. Geometric, not scalar.
%     x  GSK3beta score        y  JNK3 score
%     faint grey initial bank; matched InversionGNN rerun; immediate COMPOSE;
%     future-aware COMPOSE. Faint points for all candidates, stronger outline
%     for each nondominated frontier.
%     The panel must say "explores a stronger tradeoff surface", not "better
%     scalar score". No preference labels, no scalarization formula.
%     *** BLOCKING DEPENDENCY: this panel is NOT MEASURED YET (waits on G3).
%     It cannot be drawn from anything currently on disk. Figure 3 must not
%     ship with an invented (b). Either G3 lands first, or Figure 3 ships as
%     three panels and the caption drops the (b) clause. Do not fabricate. ***
%
% (c) RETARGETING FROM REALIZED MOLECULAR STATES.
%     MAIN VISUAL: forest / effect plot with intervals. Four effects, two
%     prefix histories each for the two main-text comparisons:
%         state reuse    +0.302 [0.186, 0.413]   +0.176 [0.069, 0.283]
%         retargeting    +0.748 [0.608, 0.891]   +1.462 [1.240, 1.691]
%     State reuse is the primary estimand. The clairvoyant arm stays in the
%     appendix and must NOT appear here.
%     INSET, small: x_0 -> x_1 -> x_2 -> x_3, then continue from x_3 under the
%     new target B, contrasted with restarting from x_0. Tiny. It only reminds
%     the reader what retargeting means; it is not the result.
%
% (d) ENDPOINT VALIDITY IS NOT PATH VALIDITY.
%     MAIN VISUAL: corridor schematic, conceptual and geometric, not
%     molecule-heavy. Shaded admissible corridor C = [2.3689, 4.4522]. One
%     trajectory that leaves the corridor and re-enters, ending endpoint-valid.
%     One trajectory held inside throughout by support restriction.
%     SMALL annotation box only: hidden-path prevalence 0.667 / 0.583,
%     Delta^V = -0.023 [-0.182, 0.134], frozen boundary -0.25.
%     This is a MECHANISM/VALIDITY panel, not a performance panel. The interval
%     straddles zero, so no empirical superiority claim may appear.
%     *** DEVELOPMENTAL: these values await the frozen 48-source confirmation.
%     The short caption below no longer carries that disclosure, so it must be
%     stated in the Sec. 5.6 prose and must remain gated by
%     \DEVELOPMENTALPATHWISE / check_submission.sh until the run lands. ***
%
% MOVED OUT OF THE FIGURE, into text or supplement: exact comparator protocols,
% the 800-source extension, additional multi-objective metrics, and the detailed
% pathwise-prevalence discussion.
%
% CAPTION ~75 words, one clause per panel.
% ==============================================================================
\begin{figure}[t]
\centering
\fbox{\parbox{0.97\textwidth}{\centering\vspace{6pt}\footnotesize
{\color{red}[Panels to be drawn. Panel (b) is not yet measured.]}\\[8pt]
\textbf{(a) Source-conditioned editing}\quad solved of $128$ vs.\ candidate allowance, measured range $k=1\ldots12$ only\\[2pt]
$31,40,44,48,51,56,61,63,65,67,68,70$, i.e.\ $54.7\%$ at $k{=}12$; comparators as horizontal markers\\[7pt]
\textbf{(b) Preference-conditioned multi-objective design}\quad JNK3 vs.\ GSK3$\beta$ objective space\\[2pt]
faint initial bank, matched rerun, immediate \method{}, future-aware \method{}, nondominated fronts \quad {\color{red}[awaits G3]}\\[7pt]
\textbf{(c) Retargeting from realized molecular states}\quad effect plot primary, molecular inset secondary\\[2pt]
state reuse $+0.302$ $[0.186,0.413]$ / $+0.176$ $[0.069,0.283]$ \quad retargeting $+0.748$ $[0.608,0.891]$ / $+1.462$ $[1.240,1.691]$\\[7pt]
\textbf{(d) Endpoint validity is not path validity}\quad prevalence and utility plotted; corridor schematic as a small supporting inset\\[2pt]
$C=[2.3689,4.4522]$; one endpoint-valid trajectory leaving and returning, beside the pathwise-restricted one\\[2pt]
prevalence $0.667$/$0.583$ \quad $\Delta^V=-0.023$ $[-0.182,0.134]$ \quad frozen boundary $-0.25$
\vspace{6pt}}}
\caption{\small\textbf{Molecular design with a frozen reference process.} \textbf{(a)}~On a $128$-source prospective panel, source-conditioned editing success increases with candidate allowance $k$. \textbf{(b)}~Preference-conditioned \method{} yields stronger tradeoffs in JNK3/GSK3$\beta$ objective space than matched baselines. \textbf{(c)}~After a goal change, continuing from the realized molecular state outperforms restarting from the source. \textbf{(d)}~Endpoint filtering does not enforce path validity, whereas support restriction makes invalid intermediate states unreachable by construction.}
\label{fig:capabilities}
\end{figure}
